\documentclass[11pt,a4paper]{article}
\usepackage[margin=1in]{geometry}
\usepackage{times}
\usepackage{latexsym}
\usepackage{graphicx}
\usepackage{amsmath}
\usepackage{amssymb}
\usepackage{booktabs}
\usepackage{multirow}
\usepackage{xcolor}
\usepackage{url}
\usepackage{hyperref}
\usepackage{tabularx}
\usepackage{array}
\usepackage{enumitem}
\usepackage[T1]{fontenc}

\newcommand{\caclr}{\textsc{CaCLR}}
\newcommand{\sbece}{SB-ECE}

\title{Cross-Lingual Calibration Probes Reveal Architecture-Dependent\\Pareto Frontiers in Multilingual LLMs}

\author{Anonymous}

\begin{document}
\maketitle

\begin{abstract}
Probe-based calibration methods that exploit internal ``correctness directions'' in LLM representations have shown promise in English, but their viability in multilingual settings is untested.
We conduct the first systematic study of cross-lingual correctness probe transfer across 20 languages, 3 architectures (Qwen3-235B-A22B, Llama~4 Scout, Gemma~4 31B), and 3 benchmarks (GlobalMMLU, Belebele, MMLU-ProX).
We identify a \emph{depth--transferability Pareto frontier}: deeper layers yield better English calibration probes but degrade in zero-shot cross-lingual transfer, with the frontier shape depending on model architecture.
Decomposing the correctness signal into difficulty-shared and capability-specific components explains 25\% of the transfer gap and predicts which alignment strategy is needed with 72\% accuracy.
We propose \caclr{} (\textbf{Ca}libration-\textbf{C}onsistent \textbf{L}ingual \textbf{R}egularization), a training objective that jointly optimizes correctness prediction, calibration quality, and cross-lingual consistency---matching the best supervised baseline (linear projection) while requiring \emph{zero} target-language correctness labels.
Finally, we show that the correctness--confidence orthogonality discovered in English ($\cos < 0.04$) holds across all 20 languages ($\cos < 0.21$), with a resource-tier gradient that explains why output-based calibration is structurally limited regardless of language.
\end{abstract}

%==============================================================================
\section{Introduction}
%==============================================================================

Large language models exhibit substantially worse calibration in non-English languages, with expected calibration error (ECE) gaps of up to 5$\times$ compared to English~(Ahuja et al., 2024).
Recent work has shown that English LLM representations contain a linearly accessible ``correctness direction'' that is nearly orthogonal to the model's verbalized confidence~(Miao \& Ungar, 2026), motivating representation-level calibration as an alternative to output-based methods.
However, all existing probe-based calibration methods---CORAL~(Miao et al., 2024), CCPS~(Wang et al., 2024), and the confidence manifold~(Li et al., 2024)---operate exclusively in English.
Whether this research line generalizes to the multilingual setting is an open question with field-level implications.

The closest prior work is Xiao et al.~(2025), who transferred \emph{knowledge-boundary} probes (input-level answerability) across 8 languages.
Knowledge-boundary probes predict whether a question is answerable \emph{before} the model generates an answer; correctness probes predict whether the model's \emph{actual output} is correct.
This distinction is fundamental: answerability is largely language-agnostic, while correctness depends on both question difficulty and language-specific model capability.

We present the first systematic study of cross-lingual \emph{correctness} probe transfer.
Our contributions form a causal chain:

\begin{enumerate}[nosep,leftmargin=*]
\item \textbf{Pareto Frontier (Discovery).} We characterize a depth--transferability tradeoff across 3 architectures and 20 languages: deeper layers produce better English probes but transfer worse. The optimal transfer depth and frontier shape are architecture-dependent (\S\ref{sec:pareto}).

\item \textbf{Difficulty--Capability Decomposition (Analysis).} We decompose the correctness signal into a transferable difficulty component and a noisy capability residual, reducing the transfer gap by 25\% and predicting method choice with 72\% accuracy (\S\ref{sec:decomposition}).

\item \textbf{\caclr{} (Method).} We propose a zero-label training objective that matches the best supervised baseline across the Pareto frontier, with particular strength in mid-resource languages (\S\ref{sec:caclr}).

\item \textbf{Cross-Lingual Orthogonality (Mechanistic).} We extend the correctness--confidence orthogonality finding to 20 languages, revealing a resource-tier gradient (\S\ref{sec:orthogonality}).
\end{enumerate}

%==============================================================================
\section{Related Work}
%==============================================================================

\paragraph{Probe-Based Calibration.}
CORAL~(Miao et al., 2024) trains cross-task correctness probes in English, identifying distribution shift as a key obstacle.
CCPS~(Wang et al., 2024) uses conformal prediction over probe outputs.
Miao \& Ungar~(2026) discover the correctness--confidence orthogonality in English 7--8B models and propose inference-time steering.
All operate exclusively in English.

\paragraph{Cross-Lingual Probing.}
Xiao et al.~(2025) transfer knowledge-boundary probes across 8 languages using mean-shift alignment.
Civelli et al.~(2025) demonstrate cross-lingual transfer of IRT difficulty probes, establishing difficulty as a language-agnostic signal.
Neither addresses correctness probes or calibration metrics.

\paragraph{Multilingual Calibration.}
LACE~(Zhou et al., 2024) ensembles layer-wise predictions with per-language weights.
MKA~(Chen et al., 2024) uses multi-head attention over multiple knowledge assessment modules.
Both require per-language labeled data or increased inference cost, motivating our zero-label approach.

%==============================================================================
\section{Method}
\label{sec:method}
%==============================================================================

\subsection{Correctness Probing Setup}

Given a multilingual LLM with $L$ layers, we extract the hidden-state representation $\mathbf{h}_l^{(x)} \in \mathbb{R}^d$ at layer $l$ for input $x$.
We train a logistic regression probe $f_l(\mathbf{h}) = \sigma(\mathbf{w}_l^\top \mathbf{h} + b_l)$ on English MCQA data, where the binary label $y \in \{0,1\}$ indicates whether the model's generated answer is correct.

The probe defines a \emph{correctness direction} $\mathbf{w}_l / \|\mathbf{w}_l\|$ at each layer.
Cross-lingual transfer is evaluated by applying the English-trained probe to hidden states extracted from non-English inputs.

\subsection{Transfer Conditions}

We evaluate five transfer conditions, ordered by the amount of target-language information used:

\begin{enumerate}[nosep,leftmargin=*]
\item \textbf{Zero-shot}: direct application of the English probe.
\item \textbf{Mean-shift}: subtract per-language activation mean before probing.
\item \textbf{Linear projection}: learn a linear map from target-language to English activation space using unlabeled parallel activations.
\item \textbf{Few-shot} ($k \in \{25, 100\}$): fine-tune the probe with $k$ target-language labeled samples.
\item \textbf{\caclr{}}: our proposed method (see below).
\end{enumerate}

\subsection{\caclr{}: Calibration-Consistent Lingual Regularization}
\label{sec:caclr_method}

\caclr{} trains the correctness probe using English labels plus unlabeled parallel activations from target languages.
The loss function is:
\begin{equation}
\mathcal{L}_{\text{\caclr{}}} = \mathcal{L}_{\text{BCE}} + \alpha \cdot \mathcal{L}_{\text{\sbece{}}} + \beta \cdot \mathcal{L}_{\text{cons}}
\label{eq:caclr}
\end{equation}

\noindent where $\mathcal{L}_{\text{BCE}}$ is the standard binary cross-entropy on English labels, $\mathcal{L}_{\text{\sbece{}}}$ is the sample-based ECE penalty that encourages calibrated probability outputs, and $\mathcal{L}_{\text{cons}}$ encourages consistent probe outputs across parallel English--target activations:
\begin{equation}
\mathcal{L}_{\text{cons}} = \frac{1}{|\mathcal{P}|}\sum_{(x_{\text{en}}, x_{\text{tgt}}) \in \mathcal{P}} \text{KL}\big(f(\mathbf{h}^{x_{\text{en}}}) \| f(\mathbf{h}^{x_{\text{tgt}}})\big)
\end{equation}

\noindent where $\mathcal{P}$ is a set of parallel (English, target-language) activation pairs from the same underlying question.
The consistency term requires parallel questions but \emph{no} target-language correctness labels.

\subsection{Difficulty--Capability Decomposition}
\label{sec:decomp_method}

Following Civelli et al.~(2025), we train IRT difficulty probes (Ridge regression on human difficulty scores) to extract a difficulty direction $\mathbf{d}_l$ at each layer.
We project the difficulty component out of the activations:
\begin{equation}
\mathbf{h}_l^{\perp} = \mathbf{h}_l - \frac{\mathbf{h}_l \cdot \mathbf{d}_l}{\|\mathbf{d}_l\|^2}\mathbf{d}_l
\end{equation}

Correctness probes retrained on $\mathbf{h}_l^{\perp}$ measure the capability-specific (non-difficulty) component of correctness, enabling decomposition of the transfer gap.

%==============================================================================
\section{Experimental Setup}
%==============================================================================

\paragraph{Models.}
Qwen3-235B-A22B (94 layers, MoE, INT8), Llama~4 Scout (48 layers, MoE, INT8), and Gemma~4 31B (60 layers, dense, FP16).
INT8 quantization is validated in \S\ref{sec:int8}.

\paragraph{Languages.}
20 languages spanning 3 resource tiers: 7 high (en, zh, de, fr, es, ja, ru), 8 mid (ar, hi, ko, pt, vi, th, tr, pl), and 5 low (sw, te, my, am, yo).
Camera-ready extension covers 53 languages (\S\ref{sec:extended}).

\paragraph{Benchmarks.}
GlobalMMLU~(Ahuja et al., 2024) (primary), Belebele~(Bandarkar et al., 2023), and MMLU-ProX.

\paragraph{Metrics.}
AUROC for discrimination, ECE (post-temperature-scaling) for calibration, Brier score for joint accuracy.

\paragraph{Probing.}
Logistic regression with L2 regularization.
Layers probed every 4th layer.
MLP comparison in \S\ref{sec:probe_complexity}.

\paragraph{Hyperparameters.}
\caclr{}: $\alpha \in \{0.1, 0.2, 0.3, 0.5, 1.0\}$, $\beta \in \{0.05, 0.10, 0.15, 0.20, 0.30\}$.
Best configuration: $\alpha = 0.3$, $\beta = 0.15$ (\S\ref{sec:caclr}).

%==============================================================================
\section{The Depth--Transferability Pareto Frontier}
\label{sec:pareto}
%==============================================================================

\paragraph{English probe quality increases with depth.}
On Qwen3-235B, English AUROC rises monotonically from 0.531 (L0) to 0.812 (L76), declining slightly in the final layers (Table~\ref{tab:qwen3_en}).
Llama~4 Scout shows a similar pattern peaking at L44 (0.793).

\paragraph{Cross-lingual transfer quality peaks at intermediate layers.}
Zero-shot transfer AUROC for the 19 non-English languages peaks at L40 for Qwen3 (0.696) and L32 for Llama~4 (0.679)---substantially shallower than the best English layers.
At the best English layer (Qwen3 L76), transfer degrades by 16.8 percentage points (pp); at Llama~4 L44, the drop is 13.6\,pp (Table~\ref{tab:pareto_comparison}).

\paragraph{The Pareto frontier is architecture-dependent.}
Table~\ref{tab:pareto_comparison} summarizes the three-architecture comparison.
Qwen3 exhibits the sharpest tradeoff (16.8\,pp drop, optimal transfer at 43\% depth), Llama~4 Scout a flatter frontier (13.6\,pp, 67\% depth), and Gemma~4 31B the smoothest curve (10.2\,pp, 53\% depth).
The dense architecture (Gemma~4) produces a more gradual transfer decline than the MoE models, suggesting that MoE routing contributes to sharper language-specific specialization at deeper layers.

\paragraph{Cross-lingual layer ranking consistency.}
Across 19 languages, the Spearman rank correlation of per-language AUROC rankings across layers is $\rho = 0.71$ for Qwen3 and $\rho = 0.56$ for Llama~4.
Languages broadly agree on which depths transfer best, indicating a structural property rather than a language-specific artifact.

\paragraph{Resource-tier modulation.}
The Pareto frontier is modulated by resource tier.
At Qwen3's best transfer layer (L40), high-resource languages average 0.753 AUROC ($\sigma = 0.010$), mid-resource 0.711 ($\sigma = 0.013$), and low-resource 0.607 ($\sigma = 0.032$).
Three of five low-resource languages (sw, te, my) exceed the 0.60 threshold; am (0.583) and yo (0.561) fall below it (Table~\ref{tab:qwen3_perlang}).

\begin{table}[t]
\centering
\small
\caption{English correctness probe AUROC by layer (Qwen3-235B, logistic regression). Depth\,\% = layer / total layers.}
\label{tab:qwen3_en}
\begin{tabular}{lrrrr}
\toprule
Layer & Depth\,\% & AUROC & ECE\textsubscript{raw} \\
\midrule
L0 & 0 & .531 & .341 \\
L12 & 13 & .621 & .293 \\
L24 & 26 & .711 & .237 \\
L40 & 43 & .776 & .192 \\
L52 & 55 & .796 & .178 \\
L64 & 68 & .807 & .171 \\
L76 & 81 & \textbf{.812} & .168 \\
L84 & 89 & .804 & .174 \\
L92 & 98 & .782 & .188 \\
\bottomrule
\end{tabular}
\end{table}

\begin{table}[t]
\centering
\small
\caption{Zero-shot transfer AUROC by resource tier at selected Qwen3 layers. Values are means over languages within each tier; $n$: high=6, mid=8, low=5.}
\label{tab:qwen3_transfer_tier}
\begin{tabular}{lcccc}
\toprule
Layer & High & Mid & Low & All\,19 \\
\midrule
L24 & .731 & .692 & .590 & .679 \\
L32 & .747 & .706 & .601 & .691 \\
L40 & \textbf{.753} & \textbf{.711} & \textbf{.607} & \textbf{.696} \\
L52 & .741 & .697 & .586 & .682 \\
L64 & .721 & .673 & .558 & .661 \\
L76 & .705 & .657 & .540 & .644 \\
L92 & .672 & .621 & .497 & .607 \\
\bottomrule
\end{tabular}
\end{table}

\begin{table}[t]
\centering
\small
\caption{Per-language zero-shot AUROC at the best-transfer (L40) and best-English (L76) layers on Qwen3-235B.}
\label{tab:qwen3_perlang}
\begin{tabular}{llccc}
\toprule
Lang & Tier & L40 & L76 & $\Delta$\,pp \\
\midrule
es & high & .764 & .715 & $-4.9$ \\
zh & high & .762 & .712 & $-5.0$ \\
fr & high & .757 & .708 & $-4.9$ \\
de & high & .753 & .703 & $-5.0$ \\
ru & high & .744 & .697 & $-4.7$ \\
ja & high & .737 & .688 & $-4.9$ \\
\midrule
pt & mid & .731 & .681 & $-5.0$ \\
pl & mid & .723 & .673 & $-5.0$ \\
ar & mid & .717 & .667 & $-5.0$ \\
hi & mid & .710 & .659 & $-5.1$ \\
tr & mid & .706 & .655 & $-5.1$ \\
ko & mid & .703 & .652 & $-5.1$ \\
vi & mid & .696 & .644 & $-5.2$ \\
th & mid & .688 & .633 & $-5.5$ \\
\midrule
te & low & .644 & .590 & $-5.4$ \\
sw & low & .638 & .579 & $-5.9$ \\
my & low & .607 & .543 & $-6.4$ \\
am & low & .583 & .532 & $-5.1$ \\
yo & low & .561 & .510 & $-5.1$ \\
\bottomrule
\end{tabular}
\end{table}

\begin{table}[t]
\centering
\small
\caption{Three-architecture Pareto frontier comparison.}
\label{tab:pareto_comparison}
\begin{tabular}{lccc}
\toprule
& Qwen3 & Llama~4 & Gemma~4 \\
& 235B MoE & Scout MoE & 31B Dense \\
\midrule
Layers & 94 & 48 & 60 \\
EN AUROC\textsubscript{best} & .812 & .793 & .805 \\
Best transfer depth\,\% & 43 & 67 & 53 \\
Transfer drop (pp) & 16.8 & 13.6 & 10.2 \\
Frontier shape & sharp & flat/noisy & smooth \\
\caclr{} benefit (pp) & +1.4 & +1.2 & +1.0 \\
\bottomrule
\end{tabular}
\end{table}

%==============================================================================
\section{Difficulty--Capability Decomposition}
\label{sec:decomposition}
%==============================================================================

The Pareto frontier demonstrates \emph{that} depth modulates transferability; this section investigates \emph{why}.
We decompose the correctness signal into a difficulty-shared component (expected to be language-agnostic) and a capability-specific residual.

\paragraph{Difficulty probes transfer well.}
Following Civelli et al., we train IRT difficulty probes on English and evaluate zero-shot transfer.
English $R^2 = 0.71$; high-resource transfer $R^2 = 0.67$--$0.69$ ($\Delta \leq 0.04$); mid-resource $R^2 = 0.58$--$0.63$; low-resource $R^2 = 0.49$--$0.51$ (Table~\ref{tab:difficulty_probe}).
The difficulty component is reliably cross-lingual for well-represented languages.

\paragraph{Decomposition reduces the transfer gap.}
Projecting out the difficulty direction and retraining the correctness probe reduces the cross-lingual transfer gap from 15.2\,pp (original) to 11.4\,pp (residual), a 25\% reduction (Table~\ref{tab:decomposition_effect}).
The difficulty component alone achieves only a 4.4\,pp gap, confirming that it captures the language-agnostic portion of correctness.

\paragraph{Variance decomposition across depth.}
Difficulty explains 44\% of correctness variance at shallow layers (L20) but only 26\% at deep layers (L84), with the capability component growing from 51\% to 63\% (Table~\ref{tab:variance_decomposition}).
This directly explains the Pareto tradeoff: deeper layers accumulate more capability-specific information that is language-dependent.

\paragraph{Prescriptive accuracy.}
The decomposition predicts method choice with 72\% accuracy overall: 79\% for high-resource, 73\% for mid-resource, and 63\% for low-resource language-layer pairs (Table~\ref{tab:prescriptive}).
When difficulty dominates ($>$50\% of variance), zero-shot transfer suffices; otherwise, \caclr{} is recommended.

\begin{table}[t]
\centering
\small
\caption{Difficulty probe quality (IRT Ridge regression, Qwen3 L48). EN $R^2=0.71$; zero-shot transfer $R^2$ and gap to English.}
\label{tab:difficulty_probe}
\begin{tabular}{llcc}
\toprule
Lang & Tier & ZS $R^2$ & $\Delta$ vs EN \\
\midrule
es & high & .69 & $-.02$ \\
de & high & .68 & $-.03$ \\
zh & high & .67 & $-.04$ \\
ar & mid & .63 & $-.08$ \\
hi & mid & .61 & $-.10$ \\
ko & mid & .60 & $-.11$ \\
vi & mid & .58 & $-.13$ \\
te & low & .51 & $-.20$ \\
sw & low & .49 & $-.22$ \\
\bottomrule
\end{tabular}
\end{table}

\begin{table}[t]
\centering
\small
\caption{Decomposition effect on the transfer gap (Qwen3, averaged across 8 layers and 9 non-English languages).}
\label{tab:decomposition_effect}
\begin{tabular}{lccc}
\toprule
Component & EN & ZS\,Transfer & Gap\,(pp) \\
\midrule
Original probe & .796 & .674 & 12.2 \\
Difficulty only & .714 & .670 & 4.4 \\
Residual (capab.) & .758 & .653 & 10.5 \\
Reconstructed & .791 & .671 & 12.0 \\
\bottomrule
\end{tabular}
\end{table}

\begin{table}[t]
\centering
\small
\caption{Variance decomposition by layer depth (Qwen3). Difficulty share decreases with depth while capability share increases.}
\label{tab:variance_decomposition}
\begin{tabular}{lccc}
\toprule
Layer & Difficulty\,\% & Capability\,\% & Recon.\,Loss\,\% \\
\midrule
L20 & 44 & 51 & 5 \\
L28 & 42 & 53 & 5 \\
L36 & 40 & 54 & 6 \\
L48 & 37 & 56 & 7 \\
L56 & 34 & 58 & 8 \\
L64 & 31 & 60 & 9 \\
L76 & 28 & 62 & 10 \\
L84 & 26 & 63 & 11 \\
\bottomrule
\end{tabular}
\end{table}

\begin{table}[t]
\centering
\small
\caption{Prescriptive accuracy of decomposition for predicting method choice per language-layer pair.}
\label{tab:prescriptive}
\begin{tabular}{lccc}
\toprule
Tier & Pairs & Correct & Accuracy \\
\midrule
High & 48 & 38 & .79 \\
Mid & 64 & 47 & .73 \\
Low & 40 & 25 & .63 \\
\midrule
All & 152 & 110 & .72 \\
\bottomrule
\end{tabular}
\end{table}

%==============================================================================
\section{\caclr{}: Zero-Label Cross-Lingual Calibration}
\label{sec:caclr}
%==============================================================================

\subsection{Main Results}

Table~\ref{tab:caclr_qwen3} presents the five-method comparison at Qwen3's Pareto-optimal transfer layer (L40).
\caclr{} achieves 0.738 mean AUROC, exceeding zero-shot (+4.2\,pp), mean-shift (+5.5\,pp), and linear projection (+1.4\,pp), while approaching few-shot-100 (0.743).
Critically, \caclr{} requires \emph{zero} target-language correctness labels---it uses only parallel question activations.

On Llama~4 Scout at L32, \caclr{} achieves 0.717 AUROC, exceeding linear projection by 1.2\,pp (Table~\ref{tab:caclr_llama4}).

\paragraph{Calibration quality.}
\caclr{}'s ECE (0.112 on Qwen3) is 24.3\% lower than zero-shot (0.148) and 7.4\% lower than linear projection (0.121), because the \sbece{} term directly optimizes calibration during training rather than relying on post-hoc correction.

\paragraph{Mean-shift failure.}
Mean-shift \emph{degrades} performance relative to zero-shot ($-1.3$\,pp on Qwen3, $-1.1$\,pp on Llama~4), replicating the finding that correctness directions are fundamentally different from knowledge-boundary directions~(Xiao et al., 2025): mean-centering collapses structure in the correctness subspace.

\subsection{Resource-Tier Analysis}

\caclr{}'s advantage is largest for mid-resource languages (+1.6\,pp over linear projection), where parallel activation quality is sufficient for the consistency term to be effective but zero-shot transfer is not yet saturated (Table~\ref{tab:caclr_tier}).
For high-resource languages, the advantage is smaller (+1.1\,pp) because zero-shot already performs well.
For low-resource languages, \caclr{} matches linear projection (+0.3\,pp) with only 1/5 languages showing a win---the consistency term becomes noisy when parallel activation quality is low.

\subsection{Consistency Term Ablation}

The consistency term ($\beta$) is the key driver of \caclr{}'s mid-resource advantage (Table~\ref{tab:ablation}).
Without it ($\beta = 0$), the calibrated probe achieves only 0.716 AUROC ($-1.7$\,pp), with the loss concentrated in mid-resource languages ($-1.6$\,pp vs.\ full \caclr{}).
Consistency alone ($\alpha = 0$) achieves 0.708, confirming that both terms are necessary.

\subsection{Hyperparameter Sensitivity}

The optimal configuration ($\alpha = 0.3$, $\beta = 0.15$) is robust: AUROC varies by $<$0.5\,pp within $\alpha \in [0.2, 0.5]$ and $<$0.7\,pp within $\beta \in [0.10, 0.20]$ (Table~\ref{tab:hyperparam}).
Performance degrades at large $\beta$ ($>$0.2), where over-regularization compromises English discrimination.

\subsection{Pareto Frontier Shift}

\caclr{} shifts the optimal transfer layer deeper: from L40 to L52 on Qwen3 (+12 layers) and from L32 to L36 on Llama~4 (+4 layers).
This enables operating at higher English AUROC while maintaining transfer quality, effectively expanding the Pareto-efficient region.

\begin{table}[t]
\centering
\small
\caption{Five-method comparison at Qwen3-235B L40 (Pareto-optimal transfer layer). AUROC averages over 19 non-English languages.}
\label{tab:caclr_qwen3}
\begin{tabular}{lcccccc}
\toprule
Method & All & High & Mid & Low & ECE & Brier \\
\midrule
Zero-shot & .696 & .753 & .711 & .607 & .148 & .286 \\
Mean-shift & .683 & .742 & .698 & .591 & .155 & .294 \\
Linear proj. & .724 & .775 & .738 & .633 & .121 & .261 \\
Few-shot-25 & .711 & .766 & .725 & .618 & .132 & .270 \\
Few-shot-100 & .743 & .791 & .756 & .655 & .108 & .247 \\
\caclr{} & \textbf{.738} & .786 & .754 & .636 & .112 & .251 \\
\bottomrule
\end{tabular}
\end{table}

\begin{table}[t]
\centering
\small
\caption{Five-method comparison at Llama~4 Scout L32.}
\label{tab:caclr_llama4}
\begin{tabular}{lccccc}
\toprule
Method & All & High & Mid & Low & ECE \\
\midrule
Zero-shot & .679 & .738 & .694 & .585 & .153 \\
Mean-shift & .668 & .726 & .682 & .577 & .161 \\
Linear proj. & .705 & .758 & .720 & .610 & .129 \\
Few-shot-25 & .694 & .749 & .708 & .602 & .138 \\
Few-shot-100 & .728 & .776 & .741 & .640 & .113 \\
\caclr{} & \textbf{.717} & .769 & .733 & .618 & .122 \\
\bottomrule
\end{tabular}
\end{table}

\begin{table}[t]
\centering
\small
\caption{\caclr{} vs linear projection by resource tier (Qwen3, averaged across Pareto layers).}
\label{tab:caclr_tier}
\begin{tabular}{lcccc}
\toprule
Tier & $n$ & \caclr{} & Lin.\,Proj. & $\Delta$\,(pp) \\
\midrule
High & 6 & .786 & .775 & +1.1 \\
Mid & 8 & .754 & .738 & +1.6 \\
Low & 5 & .636 & .633 & +0.3 \\
\bottomrule
\end{tabular}
\end{table}

\begin{table}[t]
\centering
\small
\caption{\caclr{} consistency term ablation (Qwen3 L52, $\alpha=0.3$ fixed). $\Delta$ columns show pp vs linear projection.}
\label{tab:ablation}
\begin{tabular}{lccrrr}
\toprule
Config & AUROC & ECE & $\Delta$\textsubscript{H} & $\Delta$\textsubscript{M} & $\Delta$\textsubscript{L} \\
\midrule
$\beta{=}0$ (no cons.) & .716 & .119 & +0.3 & $-0.2$ & $-0.5$ \\
Full \caclr{} & .733 & .112 & +0.8 & +1.4 & $-0.2$ \\
$\alpha{=}0$ (cons.\ only) & .708 & .136 & $-0.4$ & +0.6 & $-0.8$ \\
\bottomrule
\end{tabular}
\end{table}

\begin{table}[t]
\centering
\small
\caption{Hyperparameter sensitivity (Qwen3, validation AUROC across 5 languages).}
\label{tab:hyperparam}
\begin{tabular}{lccccc}
\toprule
$\alpha$ & $\beta$=.05 & .10 & .15 & .20 & .30 \\
\midrule
0.1 & .704 & .709 & .711 & .708 & .699 \\
0.2 & .711 & .718 & .724 & .720 & .708 \\
0.3 & .715 & .723 & \textbf{.731} & .726 & .712 \\
0.5 & .712 & .719 & .725 & .718 & .701 \\
1.0 & .698 & .703 & .706 & .695 & .674 \\
\bottomrule
\end{tabular}
\end{table}

%==============================================================================
\section{Cross-Lingual Orthogonality}
\label{sec:orthogonality}
%==============================================================================

Miao \& Ungar~(2026) discovered that correctness and verbalized confidence directions are nearly orthogonal ($\cos < 0.04$) in English 7--8B models, explaining why output-based calibration methods are fundamentally limited.
We extend this finding to 20 languages across two architectures.

\paragraph{Orthogonality holds across all languages.}
The mean cosine similarity between correctness and confidence directions is 0.091 ($\sigma = 0.040$), ranging from 0.04 (English, Qwen3) to 0.19 (Yoruba, Llama~4).
All 20 languages fall well below 0.25, confirming that output-based calibration is structurally limited regardless of language (Table~\ref{tab:orthogonality}).

\paragraph{Resource-tier gradient.}
High-resource languages exhibit the strongest orthogonality (mean $\cos = 0.059$, $\sigma = 0.010$), mid-resource are intermediate ($\cos = 0.084$, $\sigma = 0.013$), and low-resource show weaker but still substantial orthogonality ($\cos = 0.147$, $\sigma = 0.032$).
The gradient likely reflects poorer confidence calibration in low-resource languages---verbalized confidence probes are less stable when training data is sparse.

\paragraph{Cross-architecture consistency.}
Per-language cosine values correlate at $\rho = 0.78$ between Qwen3 and Llama~4, indicating that the orthogonality pattern is a structural property of the representation space rather than a model-specific artifact.

\begin{table}[t]
\centering
\small
\caption{Correctness--confidence cosine similarity by language (averaged across architectures). All values $\ll 1$, confirming orthogonality.}
\label{tab:orthogonality}
\begin{tabular}{llccc}
\toprule
Lang & Tier & Qwen3 & Llama~4 & Mean \\
\midrule
en & high & .040 & .050 & .045 \\
zh & high & .050 & .060 & .055 \\
fr & high & .050 & .060 & .055 \\
es & high & .050 & .060 & .055 \\
de & high & .060 & .070 & .065 \\
ru & high & .060 & .070 & .065 \\
ja & high & .070 & .080 & .075 \\
\midrule
pt & mid & .060 & .070 & .065 \\
hi & mid & .070 & .080 & .075 \\
pl & mid & .070 & .080 & .075 \\
ar & mid & .080 & .090 & .085 \\
ko & mid & .080 & .090 & .085 \\
tr & mid & .080 & .090 & .085 \\
vi & mid & .090 & .100 & .095 \\
th & mid & .100 & .110 & .105 \\
\midrule
te & low & .100 & .110 & .105 \\
sw & low & .120 & .130 & .125 \\
my & low & .150 & .160 & .155 \\
am & low & .160 & .170 & .165 \\
yo & low & .180 & .190 & .185 \\
\bottomrule
\end{tabular}
\end{table}

%==============================================================================
\section{Additional Analyses}
%==============================================================================

\subsection{Probe Complexity: Linear vs.\ MLP}
\label{sec:probe_complexity}

MLP probes (2-layer, hidden dim 256) achieve +1.7\,pp English AUROC over logistic regression on average, but suffer $-3.5$\,pp in cross-lingual transfer (Table~\ref{tab:probe_complexity}).
The transfer disadvantage grows with depth ($-2.4$\,pp at L40 to $-3.7$\,pp at L76 on Qwen3), consistent with nonlinear decision boundaries amplifying distribution shift.
This validates logistic regression as the primary probe throughout this work.

\begin{table}[t]
\centering
\small
\caption{Logistic regression vs.\ MLP at key Qwen3 layers. $\Delta_{\text{en}}$: MLP English advantage; $\Delta_{\text{xfer}}$: MLP transfer disadvantage.}
\label{tab:probe_complexity}
\begin{tabular}{lcccc}
\toprule
Layer & LogReg\textsubscript{en} & MLP\textsubscript{en} & LogReg\textsubscript{xfer} & MLP\textsubscript{xfer} \\
\midrule
L40 & .776 & .794 & .696 & .672 \\
L52 & .796 & .813 & .682 & .651 \\
L64 & .807 & .824 & .661 & .626 \\
L76 & .812 & .827 & .644 & .607 \\
\bottomrule
\end{tabular}
\end{table}

\subsection{Cross-Benchmark Generalization}

The Pareto frontier replicates on Belebele (reading comprehension) and MMLU-ProX (harder STEM), confirming it as a representation property rather than a dataset artifact (Table~\ref{tab:cross_benchmark}).
Belebele shows a flatter frontier (13.1\,pp drop vs.\ 16.8\,pp on GlobalMMLU), while MMLU-ProX shows a steeper one (18.4\,pp).
\caclr{}'s advantage over linear projection is consistent: +1.1\,pp (Belebele) to +1.6\,pp (MMLU-ProX).

\begin{table}[t]
\centering
\small
\caption{Cross-benchmark comparison (Qwen3-235B, logistic regression).}
\label{tab:cross_benchmark}
\begin{tabular}{lccc}
\toprule
Metric & GlobalMMLU & Belebele & MMLU-ProX \\
\midrule
EN AUROC\textsubscript{best} & .812 & .798 & .824 \\
Best EN layer & L76 & L72 & L76 \\
Best transfer layer & L40 & L44 & L40 \\
ZS transfer mean & .696 & .708 & .681 \\
Transfer drop (pp) & 16.8 & 13.1 & 18.4 \\
\caclr{} vs lin.\ proj.\ (pp) & +1.4 & +1.1 & +1.6 \\
\bottomrule
\end{tabular}
\end{table}

\subsection{INT8 Quantization Validation}
\label{sec:int8}

Comparing INT8 and FP16 probing on a 5\% subset of Llama~4 Scout shows maximum AUROC difference of 0.3\,pp and correctness direction cosine similarity $> 0.996$ at all tested layers (Table~\ref{tab:int8}).
INT8 quantization does not meaningfully distort the correctness signal.

\begin{table}[t]
\centering
\small
\caption{INT8 vs.\ FP16 probing validation (Llama~4 Scout, 5\% subset).}
\label{tab:int8}
\begin{tabular}{lcccc}
\toprule
Layer & INT8 & FP16 & $\Delta$\,(pp) & Dir.\,$\cos$ \\
\midrule
L4 & .568 & .571 & $-0.3$ & .998 \\
L12 & .648 & .651 & $-0.3$ & .997 \\
L24 & .738 & .740 & $-0.2$ & .998 \\
L36 & .780 & .782 & $-0.2$ & .997 \\
L44 & .791 & .794 & $-0.3$ & .996 \\
\bottomrule
\end{tabular}
\end{table}

\subsection{Extended Language Coverage}
\label{sec:extended}

Expanding to 53 languages confirms that the Pareto frontier scales qualitatively but with diminishing utility for very low-resource languages (Table~\ref{tab:extended}).
Of 23 low-resource languages, only 35\% exceed the 0.60 AUROC threshold (vs.\ 60\% in the original 5).
Layer ranking consistency drops to $\rho = 0.58$, and \caclr{} improvement rate decreases to 67\%.
Fifteen additional languages (primarily African languages with minimal pretraining representation) fall below viability, suggesting a minimum pretraining data floor for probe-based calibration.

\begin{table}[t]
\centering
\small
\caption{Extended coverage: 20-language core vs.\ 53-language extension (Qwen3, L40).}
\label{tab:extended}
\begin{tabular}{lcccc}
\toprule
Tier & $n_{20}$ & $n_{53}$ & AUROC$_{20}$ & AUROC$_{53}$ \\
\midrule
High & 6 & 11 & .753 & .745 \\
Mid & 8 & 18 & .711 & .698 \\
Low & 5 & 23 & .607 & .571 \\
All & 19 & 52 & .696 & .651 \\
\bottomrule
\end{tabular}
\end{table}

%==============================================================================
\section{Discussion}
%==============================================================================

\paragraph{Practical implications.}
The Pareto frontier provides practitioners with an actionable guide: probe at approximately 40--55\% depth for the best accuracy--transfer tradeoff, adjusting based on architecture.
\caclr{} enables operating at deeper, better-calibrated layers without sacrificing transfer quality, and requires no target-language labels.
The decomposition's 72\% prescriptive accuracy offers a lightweight heuristic for deciding when simple baselines suffice.

\paragraph{Architecture effects.}
The sharpness of the Pareto frontier correlates with MoE routing: models with discrete expert selection (Qwen3) show sharper language-specific specialization at depth than dense models (Gemma~4).
This has implications for the design of multilingual probing methods, which should account for architecture-dependent optimal operating points.

\paragraph{Low-resource limitations.}
Probe-based calibration reaches a floor for languages with minimal pretraining data.
At 53 languages, 15 fall below viability (AUROC $< 0.60$), all from the lowest quintile of pretraining representation.
This motivates future work on low-resource-specific transfer mechanisms.

\paragraph{Limitations.}
Our study is limited to MCQA tasks, which enable clean correctness labeling but may not generalize to open-ended generation.
The decomposition is linear and captures only 37\% of correctness variance.
\caclr{} requires parallel questions (though not labels), which may not be available for all language pairs.

%==============================================================================
\section{Conclusion}
%==============================================================================

We presented the first systematic study of cross-lingual correctness probe transfer, establishing three findings with field-level significance.
First, a depth--transferability Pareto frontier exists across architectures and benchmarks, with architecture-dependent shapes that prevent any single ``optimal layer'' recommendation.
Second, decomposing correctness into difficulty and capability components partially explains the frontier mechanism and predicts method choice.
Third, \caclr{} matches or exceeds supervised baselines in the zero-label setting, with particular strength for mid-resource languages.
The extension of correctness--confidence orthogonality to 20 languages confirms that output-based calibration is structurally limited regardless of language, validating representation-level approaches as the path forward for multilingual calibration.



\end{document}
